\par\Needspace{6\baselineskip}\medskip
\hypertarget{native-113-article084}{}\pdfbookmark[2]{84.}{article-084}
\begin{center}84.\end{center}
Since two radius vectors given in magnitude and position, together with one orbital element, suffice to determine the entire orbit, the \SourceLeaf{114}\emph{time} in which the celestial body moves from one radius vector to the other can also be determined from these data, provided that we either neglect the mass of the body or at least regard it as known. We shall adhere to the former assumption, to which the latter is easily reduced. Conversely, it is clear that two radius vectors given in magnitude and position, together with the time in which the celestial body traverses the intervening space, determine the entire orbit. This problem ranks among the most important in the theory of the motion of celestial bodies. It is rather difficult to solve, since the expression for the time in terms of the elements is transcendental and, moreover, quite complicated. This makes it all the more deserving of the utmost care. The solution to be given later seems to leave nothing further to be desired. We therefore hope that our readers will welcome our decision to save from oblivion the solution we often used before that one presented itself. It is always worthwhile to approach more difficult problems in several ways and to retain a good method even when we prefer a better one. We begin by explaining this earlier method.
